% Aspetti teorici — Slay the Roadmap / Slay the Code
\documentclass[a4paper,11pt]{report}
\usepackage[italian]{babel}
\usepackage[utf8]{inputenc}
\usepackage[T1]{fontenc}
\usepackage{geometry}
\geometry{a4paper,margin=2.5cm}
\usepackage{xcolor}
\usepackage{hyperref}
\hypersetup{colorlinks=true,linkcolor=blue!60!black,urlcolor=blue!60!black,citecolor=blue!60!black}
\usepackage{graphicx}
\usepackage{amsmath}
\usepackage{enumitem}
\usepackage{booktabs}
\graphicspath{{../../assets/}{./}}

\title{Slay the Roadmap / Slay the Code\\Aspetti Teorici della Gamification}
\author{D'Agostino (matr.\ 303226) \and Di Giacomo (matr.\ 303377)}
\date{Corso di Esempi di Gamification e Sistemi Intelligenti (EGS)\\A.A.\ 2025/2026 --- \today\\[1ex]
\small Repository: \url{https://github.com/Ago95Dev/slay_the_roadmap}}

\begin{document}
\maketitle
\tableofcontents

\chapter{Octalysis}
\label{ch:octa}
Punteggio rilevato dal radar \texttt{docs/assets/Octa\_Analysis.jpg}: \textbf{169}.

\section{Gli 8 Core Drive}
\begin{table}[h]
\centering
\begin{tabular}{@{}clr@{}}
\toprule
CD & Drive & Slay \\
\midrule
CD1 & Epic Meaning \& Calling & 5 \\
CD2 & Development \& Accomplishment & 5 \\
CD3 & Empowerment of Creativity \& Feedback & 7 \\
CD4 & Ownership \& Possession & 3 \\
CD5 & Social Influence \& Relatedness & 2 \\
CD6 & Scarcity \& Impatience & 2 \\
CD7 & Unpredictability \& Curiosity & 7 \\
CD8 & Loss \& Avoidance & 2 \\
\midrule
\multicolumn{2}{@{}l}{Totale} & 33 ($\rightarrow$ score 169 normalizzato) \\
\bottomrule
\end{tabular}
\caption{Core Drive Octalysis di Slay the Roadmap.}
\label{tab:octa}
\end{table}

Il profilo è sbilanciato su CD3 e CD7 (creatività del deck e curiosità dei
quiz/boss) con CD1/CD2 medi (missione roadmap, XP/badge) e CD4--CD6/CD8
bassi: il gioco premia l'esplorazione più che il possesso o la pressione
sociale.

\section{Mapping US $\rightarrow$ Core Drive}
\begin{itemize}[nosep]
\item US-01 (roadmap ad albero): CD1 (chiamata della missione), CD2 (progresso visibile).
\item US-02 (quiz 80\%): CD2 (mastery), CD3 (feedback immediato).
\item US-03 (ricompense/deck): CD3 (combinazioni creative), CD4 (possesso carte).
\item US-04 (boss fight): CD7 (quiz adattivi imprevedibili), CD8 (rischio sconfitta).
\item US-05 (persistenza): CD4 (proprietà del progresso), CD2 (continuità).
\end{itemize}
CD5 e CD6 restano scoperti: le leaderboard sociali e la scarsità (energia/vite)
sono mock statici (3/5 vite, 5/5 energia) e costituiscono un lavoro futuro.

\chapter{Tassonomia di Toda}
\label{ch:toda}
Il foglio ``Slay'' della tassonomia di Toda riporta 6 righe con solo 2
colonne compilate (elemento presente / esempio). Lo usiamo come metodo:
ogni riga dichiara cosa esiste oggi e cosa manca.

\begin{table}[h]
\centering\small
\begin{tabular}{@{}llp{6cm}@{}}
\toprule
Dimensione & Elemento Slay & Gap / colonne vuote \\
\midrule
Performance & Levels, Streaks & Manca curva level (level fisso 1), streak non tracciate \\
Social & Leaderboards (?) & Solo mock: nessuna classifica reale (cfr.\ Hub, Cap.~\ref{ch:hubteoria}) \\
Immersion & Avatar & Rogue Level 12 statico, nessuna personalizzazione \\
Economy & Hearts, Energy, Cards, Skills, XP, Coins & Valori fissi (vite 3/5, energia 5/5), nessuna spesa \\
Personalization & Custom path, Many choices & 1 sola roadmap navigabile su 3 \\
Feedback & Progress bar, XP notif, Correct answer feedback & Notifiche XP non collegate a HUD reale \\
\bottomrule
\end{tabular}
\caption{Le 6 righe Slay del foglio Toda con gap.}
\label{tab:toda}
\end{table}

\section{Metodo con gli esempi}
I fogli di esempio (Blank, Duolingo, EcoSteps) mostrano il procedimento:
(1) elencare gli elementi osservabili, (2) classificarli nelle 5 dimensioni
di Toda, (3) marcare esplicitamente assenze e dubbi (come il ``?'' sulle
leaderboard). Duolingo esemplifica streak/energy/campionati funzionanti;
EcoSteps mostra feedback ambientali minimali. Slay oggi sta tra i due:
ha il vocabolario (carte, cuori, energia) ma non ancora l'economia chiusa.
La Tab.~\ref{tab:toda} è quindi anche la checklist dei lavori futuri.

\chapter{Framework GamiDOC}
\label{ch:fgamidoc}
Il GamiDOC (Bassanelli) è il documento di game design richiesto tra i 5
deliverable d'esame. Le sezioni ufficiali sono:

\begin{enumerate}
\item \textbf{Context}: pubblico, bisogni, vincoli didattici (roadmap Dart/Flutter).
\item \textbf{Game behavior}: loop Studio (6--8\,min) $\rightarrow$ Quiz
(3--5 domande) $\rightarrow$ Deck (2--3 carte) $\rightarrow$ Boss (10--15\,min).
\item \textbf{Gamification}: punti, badge, level, collezioni, regole di
sblocco (soglia 80\%, tentativi illimitati ma penalizzati, energia 3+, vite 3).
\item \textbf{Aesthetics}: fantasy dungeon (Gargoyle, StackOverflow, Async,
WidgetTitan), rarità Common 6 / Rare 10 / Epic 16+Perf5 / Legendary 25+Crit25\%.
\item \textbf{Dynamics}: Dual Progression (XP giocatore + potere deck).
\item \textbf{Architettura}: rimando tecnico (si veda \texttt{specifica.pdf}).
\item \textbf{Valutazione}: protocollo utente da eseguire (si veda il
Cap.~\ref{ch:userresearch}).
\item \textbf{KPI/Analytics}: da definire (retention, mastery, tempo/boss).
\end{enumerate}

Il GamiDOC completo e presentabile è \texttt{gamidoc.pdf}; qui ne fissiamo
la struttura teorica, lì stanno i contenuti di progetto.

\chapter{User research per la Evaluation}
\label{ch:userresearch}
La User Evaluation d'esame richiede metodo, profilo, risultati e riflessioni.
Seguiamo Bassanelli: metodi \textbf{qualitativi} (think-aloud durante boss
fight e deck building, interviste semi-strutturate) + \textbf{quantitativi}
(questionari: SUS per l'usabilità, IMI per la motivazione intrinseca,
GEQ-core per l'esperienza di gioco).

\section{Protocollo (da eseguire)}
\begin{enumerate}[nosep]
\item Reclutamento: 5--8 studenti con Flutter di base (profilo da descrivere).
\item Sessione (30\,min): 5\,min onboarding, 15\,min gioco guidato
(roadmap $\rightarrow$ quiz $\rightarrow$ boss) con think-aloud,
10\,min questionari + intervista.
\item Metriche: tasso superamento quiz, tempo/boss, carte usate, SUS/IMI/GEQ.
\item Riflessioni: cosa conferma/smentisce Octalysis (CD3/CD7) e Toda (gap).
\end{enumerate}
Nessun dato utente è stato ancora raccolto: nel GamiDOC la sezione resta un
placeholder esplicito, senza inventare risultati.

\chapter{Gamification Hub}
\label{ch:hubteoria}
La vecchia istanza polyglot è spenta (HTTP 530): l'integrazione obbligatoria
è la \textbf{nuova Hub} di UniAQ.

\section{Componenti}
Console: \url{https://gamification-webapp.createlab-univaq.it};
API: \url{https://gamification-api.createlab-univaq.it/api/v1};
Swagger: \url{https://gamification-api.createlab-univaq.it/swagger-ui/index.html};
guida: \url{https://gamification-api.createlab-univaq.it/docs/api-guide.en.md};
video: \url{https://youtu.be/EJ8l2kcWiFs}.

\section{Modello concettuale}
Autenticazione \texttt{POST /auth \{username, password, origin:''GAME''\}}
$\rightarrow$ Bearer JWT valido 24\,h. Runtime
\texttt{POST /executions \{gameId, playerId, actionId, data:\{\}\}} con
\texttt{data} obbligatorio. Concetti: \textit{actions} (eventi di gioco),
\textit{point concepts} con periodi in millisecondi, \textit{badge} come
collezioni, un solo \textit{level} per point concept, \textit{regole Drools}
con \texttt{/rules/validate} e simulazioni, \textit{challenge/leaderboard/
classifiche}.

\section{Migrazione}
I vecchi \texttt{gameId} non sono validi: il gioco va ricreato da console,
le regole riscritte in Drools e validate via Swagger prima del go-live.
Questo sposta su Hub proprio ciò che Toda segnala come gap (leaderboard,
economia XP/energia) senza riscrivere il client (si veda \texttt{specifica.pdf},
Cap.~6).

\chapter{Sistemi intelligenti (Molina)}
\label{ch:molina}
Per Molina (2020) un sistema intelligente opera in un mondo complesso con
abilità cognitive --- perception, deliberation, interaction, action --- e
comportamento complesso: rationality, learning, introspection.

\section{Slay come advisor system}
Slay non è un sistema autonomo: è un \textbf{advisor system} con elementi
autonomi (quiz adattivi del boss, sblocchi automatici). La deliberation è
minima oggi (soglie fisse 25/50/75\%, danno $\mathit{round}(p/10)$), ma il
disegno prevede adattività: scelta delle domande in base agli errori passati,
difficoltà del boss in base al deck (flow), analytics di mastery per la
progressione. Questi punti collegano Molina ai debiti tecnici
(\texttt{adaptiveQuizzes} perso in \texttt{fromJson}, level fisso a 1).

\chapter{Sintesi: dalla teoria a Slay}
\label{ch:sintesi}

\begin{table}[h]
\centering\small
\begin{tabular}{@{}lp{7cm}@{}}
\toprule
Nozione & Dove è applicata in Slay \\
\midrule
CD3/CD7 Octalysis & Deck building creativo, quiz e boss imprevedibili \\
CD2 Octalysis & Quiz 80\%, badge, sblocco topic \\
CD1 Octalysis & Narrazione dungeon/roadmap \\
CD4/CD8 & Deck come possesso, rischio sconfitta boss \\
Toda Performance & Levels/streaks (parziali: level fisso) \\
Toda Economy & Hearts/Energy/Cards/XP/Coins (mock, su Hub in futuro) \\
Toda Feedback & Progress bar, feedback risposta corretta \\
GamiDOC & Loop Studio/Quiz/Deck/Boss, Dual Progression \\
Hub & Leaderboard ed economia reale (lavoro futuro) \\
Molina & Advisor system; adattività e flow come estensioni \\
\bottomrule
\end{tabular}
\caption{Tracciabilità teoria $\rightarrow$ implementazione.}
\label{tab:sintesi}
\end{table}


\begin{thebibliography}{9}
\bibitem{chou} Y.-K.\ Chou, \textit{Actionable Gamification: Beyond Points, Badges, and Leaderboards}, Octalysis Media, 2015.
\bibitem{toda} A.\ M.\ Toda et al., ``A taxonomy of game elements for gamified applications'', \textit{SIGCSE}, 2019.
\bibitem{bassanelli} S.\ Bassanelli et al., ``GamiDOC: un framework per la progettazione della gamification'', 2025 (materiale del corso EGS2025).
\bibitem{molina} M.\ Molina, ``What is an intelligent system?'', 2020 --- file \texttt{docs/legacy/What\_is\_an\_intelligent\_system.pdf}.
\bibitem{hubguide} Gamification Hub, \textit{API Guide}, \url{https://gamification-api.createlab-univaq.it/docs/api-guide.en.md}.
\bibitem{hubswagger} Gamification Hub, \textit{Swagger UI}, \url{https://gamification-api.createlab-univaq.it/swagger-ui/index.html}.
\bibitem{hubvideo} Gamification Hub, \textit{Video guida}, \url{https://youtu.be/EJ8l2kcWiFs}.
\end{thebibliography}
\end{document}
